\chapter{Conclusions}
\label{ch:conclusion}

\section{Summary}

This dissertation investigated how a complete Godot 4.6 visual behaviour-tree system can support real game-agent execution while making large trees more inspectable. The work began from two connected problems: Godot provides general graph-editor primitives rather than an integrated behaviour-tree workflow, and conventional full-detail node--link views become crowded as the number of nodes, parameters and runtime annotations grows.

The resulting version 0.9.0 plugin implements serialisable trees, typed blackboards, decorators, ordered and reactive composites, random and parallel control, repeat and wait memory, actor-method actions and conditions, reusable NPC components, undo/redo and editor-only live debugging. A complex 2D scenario demonstrates detection, pursuit, directional attack, last-known-position search, retreat and healing under tree control.

The first research question is answered positively for the defined scope: Godot's editor and resource APIs supported a practical authoring, execution and debugging pipeline. The second is answered in terms of instrumented display effects. Compact cards reduced measured card area by 55.9\%; semantic zoom reduced representative fields by 83.3\%; focus and collapse reduced visible nodes by up to 97.5\% and 99.5\%; search dimmed 99.7\% of non-target nodes; and the enhanced minimap retained 100\% node coverage. Real-GPU regression verified that these effects could be enabled and disabled without the previously observed fisheye, centring, overlap or double-edge defects. The third research question is also answered positively for the tested workload: exact-signature caching preserved execution results and reduced median tick cost by 15.98--78.00\%, with larger gains on larger trees.

\section{Contributions}

The project contributes a usable Godot plugin rather than a disconnected visual prototype; an integrated catalogue of switchable large-tree and runtime-inspection mechanisms; a safe topology-cache design; and a reproducible evidence package containing 447 passing core assertions, 21 controlled visual checks, 511 profile assertions, 53 package checks, raw measurements and fixed screenshots. It also contributes a negative but important methodological result: geometric compression must not be reported as proven human readability. The prepared participant protocol makes that boundary explicit.

\section{Limitations}

The current evidence is limited to Godot 4.6 on one Windows/GPU environment and deterministic trees no larger than 364 nodes. Automated results establish implementation properties but not long-term authoring productivity or cognitive load. The debug bridge exposes latest state rather than a full timeline. Service nodes, subtree resource reuse, Asset Library distribution and a cross-version compatibility matrix are outside the completed scope.

\section{Future work}

The immediate priority is to conduct the prepared counterbalanced study with 8--15 participants. Baseline, Compact, Collapse, Fisheye, Search and Minimap should be compared on target location, path interpretation and decorator editing. Completion time, errors, pan/zoom actions, success and subjective ratings should be analysed without merging geometrically incomparable conditions into a single unsupported score.

Further engineering work should add reusable subtree assets and explicit subtree interfaces, a debug timeline with pause/step controls, actor selection for multiple simultaneous runners, and an event-assisted topology generation counter backed by periodic exact validation. Visual work could examine animated incremental transitions and user-configurable \ac{DOI} functions while measuring revisit performance and mental-map retention. Finally, release work should validate selected Godot versions and prepare an Asset Library listing.

\section{Final statement}

The project shows that a graduation-scale implementation can move beyond imitating an existing editor. By combining executable semantics, editor-only explanation, switchable focus and complexity controls, strict regression testing and explicit evidential limits, the plugin provides both a practical AI authoring tool and a defensible platform for continued research into readable visual programming.
